% !TeX program = lualatex
% !TeX encoding = utf8
% !TeX spellcheck = uk_UA
% !TeX root = ../ElMagConspect.tex

%=========================================================
\chapter{Явище електромагнітної індукції}\label{\currfilebase}
\graphicspath{{\currfilebase/Pictures}}
%=========================================================
\epigraph{Магнітне поле не діє на нерухомі заряди --- і все ж змінне магнітне поле народжує електричне поле, здатне їх розігнати. У цьому
	факті криється, мабуть, найглибша ідея класичної електродинаміки: поле не просто \enquote{обслуговує} заряди, воно живе власним життям і саме
	породжує інші поля.
}{Переосмислення ідей сучасної фізики}



%% --------------------------------------------------------
\section{Явище електромагнітної індукції}
%% --------------------------------------------------------

%Восени 1831~р. М.~Фарадей здійснив те, що ще в 1822~р. записав собі в пам'ятній книжці як мету --- \enquoteперетворити магнетизм на електрику.
Це фундаментальне відкриття отримало назву \emphx{явища електромагнітної індукції}: воно полягає в тому, що в замкненому провідному контурі при
зміні магнітного потоку, охопленого цим контуром, виникає електричний струм, який називають \emphx{індукційним}.

%\begin{figure}[h!]\centering
%	\includegraphics[width=1\linewidth]{Faraday_in_lab}
%	\caption{Майкл Фарадей у своїй лабораторії\label{pic:FaradayLab}}
%\end{figure}

У низці дослідів Фарадей встановив, що індукційний струм виникає щоразу, коли змінюється магнітний потік крізь контур --- незалежно від \emphx{причини}
цієї зміни: чи то рухається магніт відносно нерухомого контуру, чи рухається сам контур відносно нерухомого магніту, чи змінюється сила струму в
сусідньому контурі (рис.~\ref{pic:FaradayExperiments}).

Розглянемо докладніше ці три групи дослідів.

\paragraph{Дослід 1. Рух магніту відносно нерухомої котушки.}
Котушку, з'єднану з гальванометром, закріплюють нерухомо. До неї наближають або віддаляють постійний магніт. Поки магніт нерухомий, стрілка гальванометра стоїть на нулі. Щойно магніт починає рухатися, стрілка відхиляється --- у котушці виникає індукційний струм. Під час наближення магніту струм має один напрям, під час віддалення --- протилежний. Коли рух припиняється, струм зникає.

\paragraph{Дослід 2. Рух котушки відносно нерухомого магніту.}
Магніт залишають нерухомим, а рухають котушку, з'єднану з гальванометром. Спостерігається той самий ефект: під час руху виникає індукційний струм, напрям якого залежить від напряму руху котушки. Отже, важливою є не абсолютна швидкість магніту чи котушки, а їхня \emphx{відносна} швидкість, тобто швидкість зміни магнітного потоку крізь контур.

\begin{figure}[h!]\centering
	\includegraphics[width=1\linewidth]{Faraday_experiments}
    \caption{Схема класичних дослідів Фарадея з електромагнітної індукції: виникнення індукційного струму при русі магніту відносно котушки (верхній лівий та верхній правий), при взаємному переміщенні двох котушок (нижній лівий) та при русі провідника в магнітному полі (нижній правий)\label{pic:FaradayExperiments}}
\end{figure}

\paragraph{Дослід 3. Зміна струму в сусідній котушці.}
Дві котушки розміщують поряд: одну (1) приєднують до гальванометра, другу (2) --- через ключ до джерела струму. Коли ключ замикають або розмикають, у першій котушці виникає короткочасний індукційний струм. Струм спостерігається лише в моменти \emphx{зміни} струму в другій котушці; коли струм сталий, гальванометр нічого не показує. Це явище називають \emphx{взаємоіндукцією}.

Узагальнюючи результати всіх цих дослідів, Фарадей дійшов висновку: \emphx{індукційний струм у замкненому контурі виникає тоді й лише тоді, коли
змінюється магнітний потік, охоплений цим контуром}. Сама ж причина зміни магнітного потоку --- рух магніту, рух контуру чи зміна струму в сусідньому
контурі --- на існування та напрям індукційного струму принципово не впливає.


%% --------------------------------------------------------
\section{Закон електромагнітної індукції}
%% --------------------------------------------------------

На основі дослідів Фарадея було сформульовано \emphx{закон електромагнітної індукції}: електрорушійна сила (ЕРС), що виникає в контурі,
пропорційна швидкості зміни магнітного потоку, який пронизує площу, охоплену цим контуром:
\begin{equation}\label{eq:FaradayLawInt}
	\tcbhighmath{
		\emf_\text{інд} = -\frac1c\frac{d\Phi}{dt} = -\frac1c\frac{d}{dt}\iint\limits_S \Bfield\cdot d\vect S.
	}
\end{equation}
Знак \enquote{мінус} у цьому законі --- не формальність, а прояв \emphx{правила Ленца}, з яким ми детальніше познайомимось нижче.

Магнітний потік $\Phi$, а отже, і ЕРС індукції, можуть змінюватися принципово двома різними способами:
\begin{itemize}
	\item через зміну \emphx{площі} самого контуру, що пронизується незмінним магнітним полем;
	\item через зміну самого \emphx{магнітного поля}, що пронизує нерухомий контур.
\end{itemize}
Розглянемо ці два випадки окремо --- вони мають зовсім різну фізичну природу.

%% --------------------------------------------------------
\section{Робота, що виконується при переміщенні провідника зі струмом у магнітному полі}\label{sec:Work_current}
%% --------------------------------------------------------

Візьмемо рамку, по якій може вільно ковзати перемичка, що замикає верхній і нижній проводи (рис. \ref{tikz:frame_with_rod}). Помістимо цю систему в магнітне поле. Будемо вважати спочатку, що поле перпендикулярне до площини рамки.

\begin{SCfigure}[1.25][h!]
\begin{tikzpicture}[>=latex, thick]

	% ==================== ПАРАМЕТРИ ====================
	% Основні геометричні розміри
	\pgfmathsetmacro{\railW}{6.2}         % Довжина рейок
	\pgfmathsetmacro{\railH}{4.0}         % Відстань між рейками (висота)
	\pgfmathsetmacro{\xRod}{3.5}          % Початкове положення стрижня
	\pgfmathsetmacro{\dx}{2.0}            % Зміщення dx
	\pgfmathsetmacro{\xNew}{\xRod + \dx}  % Нове положення стрижня (5.5)
	\pgfmathsetmacro{\rContact}{0.12}     % Радіус контактів

	% Крок та межі магнітного поля
	\pgfmathsetmacro{\step}{0.4}
	\pgfmathsetmacro{\xMinField}{\step}
	\pgfmathsetmacro{\xMaxField}{\railW - 0.2}
	\pgfmathsetmacro{\yMinField}{\step}
	\pgfmathsetmacro{\yMaxField}{\railH - \step}
	\pgfmathsetmacro{\xSecondField}{\xMinField + \step}
	\pgfmathsetmacro{\ySecondField}{\yMinField + \step}

	% Позиції стрілок та позначень
	\pgfmathsetmacro{\yDx}{-0.5}          % Зміщення виносної лінії dx по y
	\pgfmathsetmacro{\xJ}{\railW * 0.2}   % Позиція струму J по x
	\pgfmathsetmacro{\yF}{\railH * 0.8} % Висота вектора F по y
	% ===================================================

	% ---- Магнітне поле ----
	\foreach \x in {\xMinField, \xSecondField, ..., \xMaxField} {
		\foreach \y in {\yMinField, \ySecondField, ..., \yMaxField} {
			\fill[blue!20] (\x,\y) circle (1.2pt);
		}
	}

	% ---- Рейки та ліва сторона ----
	\draw (0,\railH) -- (\railW,\railH);
	\draw (0,0) -- (\railW,0);
	\draw (0,0) to[battery] (0,\railH);

	% ---- Стрижень ----
	\draw (\xRod,0) -- (\xRod,\railH);

	% ---- Контакти ----
	\node[circle, draw, fill=white, inner sep=2pt] (A) at (\xRod,\railH) {};
	\node[circle, draw, fill=white, inner sep=2pt] (B) at (\xRod,0) {};

	% ---- Пунктир ----
	\draw[dashed, ultra thick] (\xNew,0) -- (\xNew,\railH);
%	\draw[dashed] (\xNew,\railH) -- (\railW,\railH);
%	\draw[dashed] (\xNew,0) -- (\railW,0);


	% ---- Контакти пунктиру ----
	\node[circle, draw, fill=white, inner sep=2pt] (A1) at (\xNew,\railH) {};
	\node[circle, draw, fill=white, inner sep=2pt] (B1) at (\xNew,0) {};


	% ---- dx ----
	\draw[<->] (\xRod, \yDx) -- (\xNew, \yDx);
	\node[below] at ({(\xRod + \xNew)/2}, \yDx) {\(dx\)};

	% ---- F ----
	\draw[->, very thick, green!50! black] (\xRod, \yF) -- ++(0.9, 0) node[right] {\(\vect{F}_\text{А}\)};

	% ---- J ----
	\draw[->, very thick, red] (\xJ, 0) -- ++(0.5, 0) node[above] {\(I\)};

	% ---- B ----
	\draw[->, very thick, blue] (0.8, 1.2) -- ++(-1.1, -1.5) node[below left] {\(\Bfield\)};

    % charge
    \node[circle, draw, ball color= red!50, minimum size=0.1cm, inner sep=1,  font=\scriptsize\bfseries, text=white] (e) at (\xRod, 1) {$+$};

    \draw[->, ultra thick] (e) -- ++(1, 0) node[below] {$\vect{u}_{\perp}$};
    \draw[->, ultra thick] (e) -- ++(0, 1) node[left] {$\vect{u}_{\parallel}$};

  % Використовуємо let для обчислення кінцевої точки
    \draw[->, ultra thick, red]
        let \p1 = (e), \p2 = ($(e) + (1, 1)$) in % (1,-1) - це сума векторів v та u
        (e) -- (\p2) node[above right] {\(\vec{u}\)};

    % Для наочності - пунктирні лінії паралелограма
    \draw[dashed, gray] ($(e) + (1,0)$) -- ++(0,-1);
    \draw[dashed, gray] ($(e) + (0,1)$) -- ++(1,0);


\end{tikzpicture}
\caption{Рамка з рухомою перемичкою в магнітному полі: показано силу Ампера $\vect{F}_\text{А}$, що діє на перемичку зі струмом $I$, зміщення $dx$ та розклад повної швидкості $\vect{u}$ носія заряду на складові $\vect{u}_{\perp}$ (рух разом із перемичкою) і $\vect{u}_{\parallel}$ (рух уздовж провідника)\label{tikz:frame_with_rod}}
    \end{SCfigure}



Якщо по рамці через перемичку тече струм \(I\) (як показано на рис.~\ref{tikz:frame_with_rod}), то на перемичку діє сила Ампера:
\begin{equation}
    F_\text{А} = \frac{I}{c} \ell B, \label{eq:lorentz_force}
\end{equation}
де \(\ell\) --- довжина перемички. При зазначених на рис.~\ref{tikz:frame_with_rod} напрямках струму \(I\) і зовнішнього магнітного поля \(\Bfield\) сила \(\vect{F}\) напрямлена вправо і при зміщенні перемички на \(dx\) виконує роботу
\begin{equation}
    dA = F_\text{А} dx = \frac{I}{c} B \ell dx = \frac{I}{c} d\Phi. \label{eq:work_dx}
\end{equation}
Площа контуру зі струмом змінюється на \(dS = \ell dx\). Величина \(d\Phi = d(BS)\) являє собою зміну магнітного потоку через контур. При зміщенні перемички на кінцеву відстань
\begin{equation*}
    A_{12} = \frac{I}{c} (\Phi_2 - \Phi_1).
\end{equation*}

Нехай поле напрямлене під довільним кутом до площини контуру (рамки). Покладемо \(\Bfield = \Bfield_{\perp} + \Bfield_{\parallel}\), де \(\Bfield_{\perp}\) --- складова, напрямлена за нормаллю до площі, а \(\Bfield_{\parallel}\) --- що лежить у площині контуру. Складова \(\Bfield_{\parallel}\) створює силу, перпендикулярну до площини і, отже, перпендикулярну до зміщення перемички, і тому таку, що не виконує роботи. З урахуванням цього магнітний потік, що призводить до виконання роботи, визначається тільки компонентою поля \(\Bfield_{\perp}\):
\begin{equation*}
    \Phi = B_{\perp} S = \Bfield\vect{S}.
\end{equation*}

У випадку довільно деформованого витка і поля, що змінюється від точки до точки, весь замкнений контур потрібно подумки розбити на нескінченно малі замкнені елементи струму (рис. \ref{tikz:frame_with_rod}). Оскільки струми на спільних ділянках сусідніх витків рівні за величиною і протилежні за напрямком, то струми по всій площі вихідного контуру компенсують один одного, і залишається тільки струм \(J\) у вихідному контурі.


Вважаючи поле в межах одного елемента однорідним, отримуємо
\begin{equation*}
    dA = \frac{I}{c} d\Phi,
\end{equation*}
де магнітний потік тепер визначається формулою
\begin{equation}
    \Phi = \int\limits_{S} \Bfield d\vect{S}, \label{eq:magnetic_flux}
\end{equation}
в якій інтегрування ведеться по площі, що обмежується контуром зі струмом.

Зауважимо, що в цьому процесі виконується робота, тоді як саме магнітне поле роботу не може виконувати (сила Лоренца завжди перпендикулярна швидкості руху зарядів). Для розв'язання цієї суперечності врахуємо, що заряди в перемичці рухаються не тільки в напрямку провідника (\(\vect{u}_{\parallel}\)), а й у напрямку швидкості перемички (\(\vect{u}_{\perp}\)): \(\vect{u} = \vect{u}_{\perp} + \vect{u}_{\parallel}\). Ця друга складова швидкості призводить до появи сили Лоренца: \(\vect{F}_\text{Л} = e \vect{u}_{\perp} \times \Bfield / c\), що діє проти напрямку струму \(I\). Тому для підтримання незмінного струму в коло мають бути включені сторонні ЕРС. Саме завдяки зовнішнім джерелам енергії виконується робота з переміщення перемички.

%% --------------------------------------------------------
\section{Електромагнітна індукція в рухомих провідниках}
%% --------------------------------------------------------

%% --------------------------------------------------------
\subsection*{ЕРС індукції}
%% --------------------------------------------------------

\begin{SCfigure}[1.25][h!]
\centering
\begin{tikzpicture}[>=latex, thick]

	% ==================== ПАРАМЕТРИ ====================
	% Основні геометричні розміри
	\pgfmathsetmacro{\railW}{6.2}         % Довжина рейок
	\pgfmathsetmacro{\railH}{4.0}         % Відстань між рейками (висота)
	\pgfmathsetmacro{\xRod}{3.5}          % Початкове положення стрижня
	\pgfmathsetmacro{\dx}{2.0}            % Зміщення dx
	\pgfmathsetmacro{\xNew}{\xRod + \dx}  % Нове положення стрижня (5.5)
	\pgfmathsetmacro{\rContact}{0.12}     % Радіус контактів

	% Крок та межі магнітного поля
	\pgfmathsetmacro{\step}{0.4}
	\pgfmathsetmacro{\xMinField}{\step}
	\pgfmathsetmacro{\xMaxField}{\railW - 0.2}
	\pgfmathsetmacro{\yMinField}{\step}
	\pgfmathsetmacro{\yMaxField}{\railH - \step}
	\pgfmathsetmacro{\xSecondField}{\xMinField + \step}
	\pgfmathsetmacro{\ySecondField}{\yMinField + \step}

	% Позиції стрілок та позначень
	\pgfmathsetmacro{\yDx}{-0.5}          % Зміщення виносної лінії dx по y
	\pgfmathsetmacro{\xJ}{\railW * 0.4}   % Позиція струму J по x
	\pgfmathsetmacro{\yF}{\railH * 0.6} % Висота вектора F по y
	% ===================================================

	% ---- Магнітне поле ----
	\foreach \x in {\xMinField, \xSecondField, ..., \xMaxField} {
		\foreach \y in {\yMinField, \ySecondField, ..., \yMaxField} {
			\fill[blue!40] (\x,\y) circle (1.2pt);
		}
	}

	% ---- Рейки та ліва сторона ----
	\draw (0,\railH) -- (\railW,\railH);
	\draw (0,0) -- (\railW,0);
	\draw (0,0) -- (0,\railH);

	% ---- Стрижень ----
	\draw (\xRod,0) -- (\xRod,\railH);

	% ---- Контакти пунктиру ----
	\node[circle, draw, fill=white, inner sep=2pt] (A1) at (\xNew,\railH) {};
	\node[circle, draw, fill=white, inner sep=2pt] (B1) at (\xNew,0) {};

	% ---- Пунктир ----
	\draw[dashed, ultra thick] (\xNew,0) -- (\xNew,\railH);
%	\draw[dashed] (\xNew,\railH) -- (\railW,\railH);
%	\draw[dashed] (\xNew,0) -- (\railW,0);


	% ---- Контакти пунктиру ----
	\node[circle, draw, fill=white, inner sep=2pt] (A1) at (\xNew,\railH) {};
	\node[circle, draw, fill=white, inner sep=2pt] (B1) at (\xNew,0) {};


	% ---- dx ----
	\draw[<->] (\xRod, \yDx) -- (\xNew, \yDx);
	\node[below] at ({(\xRod + \xNew)/2}, \yDx) {\(dx\)};

	% ---- F ----
%	\draw[->, very thick, green!50! black] (\xRod, \yF) -- ++(0.9, 0) node[right] {\(\vect{F}\)};

	% ---- J ----
	\draw[->, very thick, red] (\xJ, 0) -- ++(-0.5, 0) node[above] {\(I\)};

	% ---- B ----
	\draw[->, very thick, blue] ({\xRod/2},{\railH/2}) -- ++(-1.1, -1.5) node[below left] {\(\Bfield\)};

    % charge
    \node[circle, draw, ball color= red!50, minimum size=0.1cm, inner sep=1,  font=\scriptsize\bfseries, text=white] (e) at (\xRod, 2.5) {$+$};

    \draw[->, ultra thick] (e) -- ++(1, 0) node[below] {$\vec{v}$};
    \draw[->, ultra thick] (e) -- ++(0, -1) node[right] {$\vec{u}$};

    \draw[thick, ->, red] ({\xRod/2},{\railH/2}) [partial ellipse=45:{270+45}:1];

\end{tikzpicture}
\caption{Виникнення ЕРС індукції в рухомій перемичці: на позитивний заряд $+$ у провіднику, що рухається зі швидкістю $\vect{v}$, діє сила Лоренца $\vect{F}_\text{Л}$, напрямлена вздовж провідника (швидкість $\vect{u}$). Ця сила викликає індукційний струм $I$ у контурі\label{tikz:move_rod}}
    \end{SCfigure}

Розглянемо провідну рамку, замкнену рухомою перемичкою (рис.~\ref{tikz:move_rod}). Нехай рамка поміщена в магнітне поле з індукцією \(\Bfield\), напрямленою перпендикулярно до площини. Якщо перемичка рухається зі швидкістю \(\vect{v}\), то на заряди в перемичці діє сила Лоренца:
\begin{equation*}
    \vect{F}_\text{Л} = \frac{q}{c} \vect{v} \times \Bfield.
\end{equation*}
Як видно з рис. \ref{tikz:move_rod}, ця сила приводить заряди в рух, викликаючи індукційний струм у від'ємному напрямку обходу контуру. Вона еквівалентна силі, що створюється стороннім електричним полем
\begin{equation*}
    \Efield_{\text{стор}} = \frac{\vect{F}}{q} = \frac{1}{c} \vect{v} \times \Bfield
\end{equation*}
і викликає в перемичці довжиною \(l\) ЕРС індукції, рівну
\begin{equation*}
    \emf_{\text{інд}} = \oint \Efield_{\text{стор}} d\vect{l} = - E_{\text{стор}} l = - \frac{1}{c} B \ell v.
\end{equation*}
Знак «\(-\)» тут пов'язаний з тим, що \(\mathcal{E}_{\text{інд}}\) створює струм у від'ємному напрямку (додатний напрямок задається вектором \(\Bfield\)).

%\placeholderfig{Зліва --- рамка зі струмом у магнітному полі. Перемичка може вільно ковзати, безперервно замикаючи верхній і нижній проводи рамки; справа --- напрямок сили Лоренца, що діє на рухомі позитивні заряди}{fig:5.2.1}

Оскільки \(\ell v = dS/dt\), де \(S\) --- площа контуру, то
\begin{equation}\label{eq:faraday_law_for_moving}
    \emf_{\text{інд}} = - \frac{1}{c} \frac{d\Phi}{dt},
\end{equation}
де \(\Phi = BS\) --- магнітний потік через контур.

Якщо є складова магнітного поля \(\Bfield_{\parallel}\), паралельна площині контуру, то вона не призводить до появи індукційного струму вздовж контуру, оскільки відповідна складова сили Лоренца \(\vect{F}^{(1)} = \frac{1}{c} \vect{v} \times \Bfield_{\parallel}\) перпендикулярна до площини контуру, тобто напрямлена перпендикулярно до провідника. Це означає, що у виразі для потоку \(\Phi\) потрібно врахувати тільки складову \(\Bfield_{\perp}\), тобто покласти \(\Phi = B_{\perp} S = \Bfield\vect{S}\).

Якщо магнітне поле неоднорідне в границях контуру, то потрібно розбити контур на нескінченно малі замкнені елементи струму і розглянути рух цих елементів, вважаючи поле однорідним у межах окремого елемента. Отже, магнітний потік визначається інтегралом по площі, що обмежується контуром:
\begin{equation*}
    \Phi = \int\limits_{S} \Bfield d\vect{S}
\end{equation*}


\begin{Attention}
Отже, у розглянутому випадку \emphx{безпосередньою фізичною причиною виникнення ЕРС індукції є сила Лоренца}, що діє на вільні носії заряду в рухомій перемичці. Саме вона «штовхає» електрони вздовж провідника, створюючи індукційний струм.

Однак тут виникає природне запитання: а що буде, якщо перейти в систему відліку, пов'язану з самою перемичкою? Адже відносно неї провідник нерухомий, і магнітна складова сили Лоренца, яка залежить від швидкості, дорівнюватиме нулю. Якщо причина зникла, то чому тоді в контурі все одно тече струм?

Відповідь на це запитання виявляється напрочуд глибокою і вимагає перегляду наших уявлень про електричне та магнітне поля. Виявляється, те, що ми називаємо «причиною» виникнення ЕРС, залежить від вибору системи відліку. У системі, де провідник рухається, причина --- сила Лоренца. Але в системі, де провідник нерухомий, причина буде зовсім іншою: в ній виникне \emphx{вихрове електричне поле}, породжене змінним магнітним полем.

Докладно цей механізм ми розглянемо в наступному розділі, присвяченому ЕРС індукції в \emphx{нерухомих} провідниках. А питання про те, чому обидва описи --- і через силу Лоренца, і через вихрове електричне поле --- дають однаковий результат, тісно пов'язане з принципом відносності та буде обговорене пізніше.
\end{Attention}

%% --------------------------------------------------------
\subsection*{Принцип роботи генераторів змінного струму}
%% --------------------------------------------------------

Явище електромагнітної індукції в рухомих провідниках лежить в основі роботи \emphx{генераторів електричного струму} --- пристроїв, без яких
неможливо уявити сучасну електроенергетику: у них механічна енергія обертання перетворюється на електричну.

\begin{figure}[h!]\centering
	\includegraphics[width=0.45\linewidth]{AC-Generator}
	\caption{Схема найпростішого генератора змінного струму\label{pic:ACGenerator}}
\end{figure}

Найпростіший генератор складається з рамки, кінці якої закріплені на кільцях, що обертаються разом з нею; до кілець щільно прилягають нерухомі
щітки-контакти (рис.~\ref{pic:ACGenerator}). Якщо рамку площею $S$ обертати в однорідному полі $\Bfield$ з кутовою швидкістю $\omega$, то потік крізь
неї
\begin{equation*}
	\Phi(t) = BS\cos(\omega t), \qquad \frac{d\Phi}{dt} = -BS\omega\sin(\omega t),
\end{equation*}
тому в рамці виникає ЕРС індукції
\begin{equation}\label{eq:GeneratorEMF}
	\tcbhighmath{
		\emf_\text{інд} = \frac{BS\omega}{c}\sin(\omega t),
	}
\end{equation}
а якщо на щітки підключити навантаження опором $R$, у колі потече \emphx{змінний} індукційний струм
\begin{equation}\label{eq:GeneratorCurrent}
	I(t) = \frac{\emf_\text{інд}}{R} = \frac{BS\omega}{cR}\sin(\omega t),
\end{equation}
що змінюється за законом синуса з частотою $\omega$ --- частотою обертання рамки. Саме такий принцип лежить в основі роботи практично всіх
електростанцій, що генерують змінний струм.

%% --------------------------------------------------------
\section{Електромагнітна індукція в нерухомих провідниках}
%% --------------------------------------------------------

%% --------------------------------------------------------
\subsection*{Виникнення ЕРС за рахунок зміни магнітного поля}
%% --------------------------------------------------------

Розглянемо тепер принципово інший випадок: контур нерухомий і незмінний за формою, а змінюється саме магнітне поле, що його пронизує
(рис.~\ref{pic:FieldChange}):
\begin{equation}\label{eq:FaradayFieldChange}
	\emf_\text{інд} = -\frac1c\frac{d\Phi}{dt} = -\frac1c\parttime{\Bfield}\cdot\vect S.
\end{equation}

\begin{figure}[h!]\centering
	\includegraphics[width=0.5\linewidth]{Field_change}
	\caption{ЕРС індукції в нерухомому контурі за рахунок зміни поля $\Bfield$\label{pic:FieldChange}}
\end{figure}

Тут виникає принципове запитання: \emphx{яка причина виникнення ЕРС}? Носії заряду в нерухомому провіднику мають нульову швидкість, а магнітне
поле, я\emphx{не діє на нерухомі заряди} взагалі! То за рахунок чого ж вони прискорюються? Відповідь
на це запитання --- \emphx{вихрове електричне поле} --- ми дамо трохи згодом; спершу розглянемо кілька важливих якісних наслідків закону
індукції.

%% --------------------------------------------------------
\subsection*{Правило Ленца}
%% --------------------------------------------------------

\emphx{Правило Ленца} визначає \emphx{напрям} індукційного струму: індукційний струм завжди має такий напрям, щоб своїм власним магнітним полем
\emphx{протидіяти} зміні магнітного потоку, яка його породжує, тобто щоб послабити причину, яка збуджує цей струм (рис.~\ref{pic:LenzRule}).

\begin{figure}[h!]\centering
	\localinput{lenz-rule.tikz}
	\caption{Правило Ленца: напрям індукційного струму залежить від знака $\partial\Bfield/\partial
		t$\label{pic:LenzRule}}
\end{figure}

\begin{Attention}
	Правило Ленца --- не окремий незалежний закон природи, а прямий наслідок знака \enquote{мінус} у законі
	індукції~\eqref{eq:FaradayLawInt}, який, своєю чергою, диктується законом збереження енергії: якби індукційний струм \emphx{підсилював} зміну
	потоку (а не протидіяв їй), процес наростав би лавиноподібно й необмежено, породжуючи енергію з нічого. Це один із проявів загальнішого
	\emphx{принципу Ле Шательє--Брауна}: будь-який зовнішній вплив, що намагається вивести систему зі стійкої рівноваги, породжує в ній
	компенсаційні процеси, спрямовані на послаблення цього впливу.
\end{Attention}

%% --------------------------------------------------------
\subsection*{Струми Фуко}
%% --------------------------------------------------------

\emphx{Струмами Фуко} (вихровими струмами) називають індукційні струми, що виникають не в тонких лінійних провідниках, а в товщі суцільних
масивних провідників при зміні магнітного потоку крізь них (рис.~\ref{pic:FukoCurrents}). Вони підпорядковані тому самому правилу Ленца: їхнє
магнітне поле напрямлене так, щоб протидіяти зміні потоку, яка їх породжує.

\begin{SCfigure}[1][h!]\centering
	\includegraphics[width=0.45\linewidth]{Eddy_currents}
	\caption{Індукування вихрових струмів Фуко в провідному диску, що обертається в полі магніту\label{pic:FukoCurrents}}
\end{SCfigure}

\begin{Attention}
	Оскільки струми Фуко течуть у товщі провідника (а не по тонкому дроту з відомим опором), вони можуть спричиняти \emphx{значне} тепловиділення
	(закон Джоуля--Ленца, гл.~\ref{SteadyCurrent}) --- саме на цьому явищі, зокрема, ґрунтується \emphx{індукційне нагрівання} металів. З
	іншого боку, ті самі струми, взаємодіючи з магнітним полем (сила Ампера), створюють силу, що \emphx{гальмує} рух провідника відносно джерела
	поля --- явище \emphx{магнітного демпфування}, яке використовують, наприклад, в електромагнітних гальмах і в заспокоювачах коливань
	вимірювальних приладів.
\end{Attention}

%% --------------------------------------------------------
\subsection*{Вихрове електричне поле}
%% --------------------------------------------------------

Повернімося до питання, яке ми залишили відкритим: що саме прискорює нерухомі заряди при зміні $\Bfield$? Оскільки $\Phi = \iint_S \Bfield\cdot
d\vect S$, а ЕРС, за означенням~\eqref{eq:EMFDef} (гл.~\ref{SteadyCurrent}), дорівнює $\emf = \oint_L \Efield\cdot d\vect\ell$, із закону
індукції безпосередньо випливає:
\begin{equation}\label{eq:CurlEInt}
	\tcbhighmath{
		\oint\limits_L \Efield\cdot d\vect\ell = -\frac1c\iint\limits_S \frac{\partial\Bfield}{\partial t}\cdot d\vect S.
	}
\end{equation}
Це --- \emphx{теорема про циркуляцію} електричного поля: циркуляція напруженості електричного поля по довільному контуру дорівнює, з протилежним
знаком, потокові швидкості зміни індукції магнітного поля крізь довільну поверхню, натягнуту на цей контур. За теоремою Стокса, цю теорему можна
записати й у диференціальній формі:
\begin{equation}\label{eq:CurlEDiff}
	\tcbhighmath{
		\rot\Efield = -\frac1c\frac{\partial\Bfield}{\partial t}.
	}
\end{equation}

За Максвеллом, явище електромагнітної індукції полягає саме в тому, що \emphx{будь-яке змінне магнітне поле збуджує в просторі електричне поле};
провідники для цього взагалі не потрібні --- вони лише дають змогу побачити це поле через індукційний струм, який воно викликає в них
(рис.~\ref{pic:VortexField}).

\begin{SCfigure}[1.5][h!]\centering
	\localinput{vortex-field.tikz}
	\caption{Вихрове електричне поле, породжене зростаючим у часі магнітним полем\label{pic:VortexField}}
\end{SCfigure}

\begin{Attention}
	На відміну від електростатики, де завжди $\rot\Efield=0$ (поле потенціальне), у присутності змінного магнітного поля $\rot\Efield\neq0$. Це
	принципово важлива відмінність: індуковане електричне поле не є потенціальним --- воно \emphx{вихрове}: його силові лінії ніде не починаються
	й не закінчуються (на відміну від ліній електростатичного поля, які завжди починаються й закінчуються на зарядах), а замкнені самі на себе,
	подібно до силових ліній магнітного поля. Електричне поле, породжене зміною магнітного поля, має \emphx{зовсім іншу структуру}, ніж
	електростатичне, і не пов'язане безпосередньо з жодними електричними зарядами.

	Чому ж тоді це поле взагалі називають \emphx{електричним}, якщо воно має інше походження й іншу конфігурацію ліній, ніж звичне статичне поле?
	Відповідь проста й фундаментальна: \emphx{вихрове поле діє на заряд точнісінько так само, як і електростатичне}, --- а саме це ми з самого
	початку курсу вважали \emphx{головною визначальною властивістю} електричного поля. Саме на прискоренні заряджених частинок вихровим електричним
	полем ґрунтується робота \emphx{бетатрона} --- прискорювача електронів, у якому змінне поле потужного електромагніту породжує вихори
	електричного поля, здатні розганяти електрони до швидкостей, близьких до швидкості світла, причому електричний струм виникає безпосередньо у
	вакуумній камері, без жодних металевих провідників.
\end{Attention}

%% --------------------------------------------------------
\section{Явище самоіндукції}
%% --------------------------------------------------------

%% --------------------------------------------------------
\subsection*{Індуктивність контуру}
%% --------------------------------------------------------

Зміна струму в довільному контурі змінює й створюване ним магнітне поле, а отже --- і власний магнітний потік крізь цей самий контур; за законом
індукції це, своєю чергою, породжує в контурі ЕРС. Це явище називають \emphx{самоіндукцією} (рис.~\ref{pic:SelfInduction}).

\begin{SCfigure}[1.25][h!]\centering
	\localinput{self-induction-field.tikz}
	\caption{Контур зі струмом $I$ і створюване ним власне магнітне поле\label{pic:SelfInduction}}
\end{SCfigure}

Якщо поблизу контуру немає феромагнетиків, поле $\Bfield$, а отже, й повний магнітний потік $\Phi$ крізь контур пропорційні силі струму $I$ в
ньому:
\begin{equation}\label{eq:SelfInductionFlux}
	\tcbhighmath{
		\Phi = \frac1c LI,
	}
\end{equation}
де коефіцієнт пропорційності $L$ називають \emphx{індуктивністю контуру} --- суто геометричною характеристикою (залежною від форми, розмірів
контуру й від магнітних властивостей середовища навколо нього), яка не залежить від сили струму. При зміні струму в контурі, за законом Фарадея,
виникає \emphx{ЕРС самоіндукції}:
\begin{equation}\label{eq:SelfInductionEMF}
	\tcbhighmath{
		\emf_\text{сі} = -\frac{L}{c^2}\frac{dI}{dt}.
	}
\end{equation}

\begin{Attention}
	Знак \enquote{мінус} тут --- знову прояв правила Ленца: ЕРС самоіндукції завжди напрямлена так, щоб протидіяти зміні сили струму в контурі ---
	вона гальмує зростання струму й \emphx{підтримує} його спадання. Через цю властивість контур з великою індуктивністю поводиться подібно до
	тіла з великою механічною інерцією (масою): він \enquote{опирається} будь-якій швидкій зміні струму.
\end{Attention}

%% --------------------------------------------------------
\subsection*{Приклади обчислення індуктивності}
%% --------------------------------------------------------

\emphx{Соленоїд.} Магнітне поле всередині нескінченного соленоїда з магнетиком проникності $\mu$, густиною намотки $n=N/l$ витків на одиницю
довжини (гл.~\ref{MagFieldMedia}), дорівнює $B = 4\pi\mu NI/(cl)$. Повний магнітний потік крізь усі $N$ витків соленоїда $\Phi = NBS =
\dfrac{4\pi\mu N^2}{cl}IS$; порівнюючи з означенням~\eqref{eq:SelfInductionFlux}, отримуємо індуктивність соленоїда:
\begin{equation}\label{eq:SolenoidInductance}
	\tcbhighmath{
		L = \frac{4\pi\mu N^2 S}{l}.
	}
\end{equation}

\emphx{Тороїд.} Для тороїдальної котушки з $N$ витками, прямокутним перерізом ширини $b-a$ (де $a$ і $b$ --- внутрішній і зовнішній радіуси) і
висотою $d$ (рис.~\ref{pic:Toroid}), аналогічний (дещо громіздкіший) розрахунок дає:

\begin{figure}[h!]\centering
	\includegraphics[width=0.3\linewidth]{toroid}
	\caption{Тороїдальна котушка\label{pic:Toroid}}
\end{figure}
\begin{equation}\label{eq:ToroidInductance}
	L = 2N^2 d\ln\left(\frac{b}{a}\right).
\end{equation}

\begin{Attention}
	В гауссовій системі одиниць індуктивність, як неважко переконатись з формул~\eqref{eq:SolenoidInductance}--\eqref{eq:ToroidInductance},
	вимірюється просто в \emphx{сантиметрах}, $[L]=\text{см}$ --- ще одне яскраве підтвердження того, що індуктивність, як і ємність
	(гл.~\ref{ElEnergy}), --- суто \emphx{геометрична} характеристика контуру.
\end{Attention}

%% --------------------------------------------------------
\subsection*{Перехідні процеси в колі з індуктивністю}
%% --------------------------------------------------------

\emphx{Встановлення струму в $LR$-колі.} Розглянемо коло з послідовно з'єднаних джерела ЕРС $\emf$, котушки індуктивності $L$, резистора $R$
і ключа $S$ (рис.~\ref{pic:LRTurnOn}). Після замикання ключа закон Ома для кола, $\emf + \emf_\text{сі} = IR$, з урахуванням
формули~\eqref{eq:SelfInductionEMF}, набуває вигляду
\begin{equation}\label{eq:LRTurnOnEq}
	\frac{L}{c^2}\frac{dI}{dt} + IR = \emf.
\end{equation}

\begin{figure}[h!]\centering
	\localinput{lr-turnon.tikz}
	\caption{Встановлення струму в колі з індуктивністю після замикання ключа\label{pic:LRTurnOn}}
\end{figure}

Це рівняння --- повний аналог рівняння зарядки конденсатора (гл.~\ref{SteadyCurrent}), і його розв'язок (з початковою умовою $I(0)=0$) --- струм,
що асимптотично наближається до стаціонарного значення $\emf/R$:
\begin{equation}\label{eq:LRTurnOnSolution}
	\tcbhighmath{
		I(t) = \frac{\emf}{R}\left(1-e^{-t/\tau}\right), \qquad \tau = \frac{L}{c^2R}
	}
\end{equation}
--- сталу часу $\tau$ називають \emphx{часом релаксації} $LR$-кола.

\emphx{Екстраструми при розмиканні.} Нехай тепер котушка з індуктивністю $L$ і власним (внутрішнім) опором обмотки $r$ спочатку під'єднана до
джерела ЕРС $\emf$ через ключ $S$ і резистор $R$, під'єднаний паралельно котушці (рис.~\ref{pic:LRTurnOff}); у сталому режимі через котушку
тече струм $I_0 = \emf/r$.

\begin{figure}[h!]\centering
	\localinput{lr-turnoff.tikz}
	\caption{Екстраструм розмикання в колі з індуктивністю\label{pic:LRTurnOff}}
\end{figure}

Після розмикання ключа $S$ (відключення джерела) магнітне поле котушки, а отже, і потік крізь неї, починають спадати. Це, за законом Ленца,
збуджує ЕРС самоіндукції $\emf_\text{сі}$ і підтримуваний нею \emphx{екстраструм розмикання} --- індукційний струм, що тепер тече вже через
резистор $R$ (замкнене коло $L$--$r$--$R$ без джерела). Закон Ома для цього контуру: $I(R+r) = -\dfrac{L}{c^2}\dfrac{dI}{dt}$, звідки
\begin{equation}\label{eq:LRTurnOffSolution}
	\tcbhighmath{
		I(t) = I_0\, e^{-(R+r)t/(Lc^{-2})} = I_0\, e^{-t(R+r)c^2/L}.
	}
\end{equation}

\begin{Attention}
	Якщо $R\gg r$ (типова ситуація, коли навантаження має значно більший опір, ніж провід самої котушки), то в момент розмикання ЕРС самоіндукції
	досягає значення $\emf_\text{сі} = \dfrac{R}{r}\emf\, e^{-t(R+r)c^2/L}$, яке при $R\gg r$ може \emphx{значно перевищувати} ЕРС самого
	джерела $\emf$! Саме тому раптове розмикання кола з великою індуктивністю (наприклад, вимкнення потужного електромагніту чи реле) часто
	супроводжується електричним пробоєм або яскравою іскрою на контактах ключа --- індуктивність \enquote{чинить опір} різкому зникненню струму
	настільки завзято, що готова створити на короткий час дуже високу напругу.
\end{Attention}

%% --------------------------------------------------------
\section*{Підсумки}
%% --------------------------------------------------------

Явище електромагнітної індукції є тією вирішальною ланкою, що остаточно об'єднує електрику та магнетизм у єдину цілісну теорію: змінне магнітне поле генерує вихрове електричне, а сила Лоренца (розд.~\ref{MagFieldVacuum}) демонструє фундаментальний зв'язок магнітного поля з рухом зарядів.

Явище самоіндукції свідчить про те, що встановлення струму в контурі вимагає виконання роботи проти індукованого електричного поля. Наступний логічний крок --- дослідити, як ця робота акумулюється в просторі, розрахувати \emphx{енергію магнітного поля}.

\begin{Questions}
\begin{quizlist}
	\item Сформулюйте суть явища електромагнітної індукції. Опишіть три групи дослідів Фарадея (рис.~\ref{pic:FaradayExperiments}). Що в них спільного і чому важливою є лише відносна швидкість магніту та котушки?
	\item Коли в контурі виникає індукційний струм, а коли ні? Чому сама причина зміни магнітного потоку не впливає на існування струму?
	\item Запишіть закон електромагнітної індукції в інтегральній формі. Що означає знак «мінус» і якими двома способами може змінюватися магнітний потік крізь контур?
	\item Виведіть вираз для роботи сили Ампера при зміщенні перемички на $dx$ (рис.~\ref{tikz:frame_with_rod}) і покажіть, що $dA = \frac{I}{c}\,d\Phi$. Яка складова поля $\Bfield$ виконує роботу, а яка ні?
	\item Магнітне поле не виконує роботи, але при переміщенні перемички зі струмом робота виконується. Як розв'язується ця суперечність? Що таке складові швидкості $\vect{u}_{\perp}$ та $\vect{u}_{\parallel}$ і яку роль відіграють сторонні ЕРС?
	\item Що є безпосередньою причиною ЕРС індукції в рухомій перемичці (рис.~\ref{tikz:move_rod})? Виведіть вираз $\emf_\text{інд} = -B\ell v/c$ і покажіть, що він дає закон Фарадея.
	\item Чому в рухомому провіднику причиною ЕРС є сила Лоренца, а в системі відліку, пов'язаній із самим провідником, --- вихрове електричне поле? Про що це свідчить щодо співвідношення електричного й магнітного полів?
	\item Поясніть принцип роботи найпростішого генератора змінного струму (рис.~\ref{pic:ACGenerator}). Виведіть формулу для ЕРС та струму в рамці, що обертається в однорідному полі.
	\item Яка причина ЕРС індукції в нерухомому контурі, якщо магнітне поле не діє на нерухомі заряди?
	\item Сформулюйте правило Ленца (рис.~\ref{pic:LenzRule}). Чому воно є наслідком закону збереження енергії? Що сталося б, якби індукційний струм підсилював зміну потоку?
	\item Що таке струми Фуко? Наведіть приклади їх корисного (індукційне нагрівання, магнітне демпфування) та шкідливого прояву.
	\item Запишіть теорему про циркуляцію електричного поля в інтегральній та диференціальній формі. У чому принципова відмінність вихрового електричного поля від електростатичного? Чому його все ж називають електричним?
	\item Що таке самоіндукція та індуктивність контуру? Від чого вона залежить, а від чого ні? Запишіть ЕРС самоіндукції та поясніть її інерційну роль у колі.
	\item Виведіть індуктивність довгого соленоїда. У яких одиницях вимірюється індуктивність у гаусовій системі і чому це підтверджує її геометричний характер?
	\item Виведіть закон зростання струму в $LR$-колі після замикання ключа. Що таке час релаксації $\tau$ і від чого він залежить?
	\item Що таке екстраструм розмикання? Чому при $R\gg r$ ЕРС самоіндукції в момент розмикання може значно перевищувати ЕРС джерела і до яких наслідків це призводить?
\end{quizlist}
\end{Questions}